% Copyright (c) 2026 Eugene Joseph M. Ragasa.
% SPDX-License-Identifier: Apache-2.0
% Author: Eugene Joseph M. Ragasa
% Formal proof completed: 2026-10-19
\documentclass[11pt]{article}
\usepackage[T1]{fontenc}\usepackage{amsmath,amssymb}\usepackage[margin=1in]{geometry}\usepackage[hidelinks]{hyperref}
\title{Proof Dossier: Real-Rotation Action on Pauli Coordinates}\author{Eugene Joseph M. Ragasa}\date{Formal proof completed: 2026-10-19}
\begin{document}\maketitle
\section{Identity and status}
ID \texttt{BG-FOUNDATION-THM-003}; status \texttt{proved}. Lean declaration:
\par\noindent {\scriptsize\path{ProjectKoios.BandGauge.Foundations.rankTwoRealRotation_conjTranspose_mul_pauli_mul}}.
\section{Statement and proof}
For $D=x\sigma_x+y\sigma_y+z\sigma_z$,
\[R(\phi)^\dagger D R(\phi)=
 (x\cos2\phi-z\sin2\phi)\sigma_x+y\sigma_y+
 (x\sin2\phi+z\cos2\phi)\sigma_z.\]
The proof is an explicit $2$-by-$2$ matrix calculation using the double-angle formulas.
\section{Boundary}
This is standard supporting matrix algebra, not a paper novelty claim. Kernel-reported
axioms: \texttt{propext}, \texttt{Classical.choice}, \texttt{Quot.sound}. No admissions
or scientific axioms. AI-assisted and kernel-checked; no independent expert review recorded.
\end{document}
